\chapter{Ing\'enierie de variables (feature engineering)}
\label{ch:feature-engineering}

\begin{quote}
\emph{<< Trouver de bonnes variables est difficile, chronophage, demande de l'expertise. Le machine learning appliqu\'e, c'est essentiellement du feature engineering. >>}
--- Andrew Ng
\end{quote}

% ============================================================
\section{Qu'est-ce que le feature engineering ?}

Le feature engineering est le processus de cr\'eation de nouvelles variables d'entr\'ee \`a partir des donn\'ees existantes pour am\'eliorer la performance du mod\`ele. Alors que l'encodage et la mise \`a l'\'echelle (Chapitre~5) transforment les variables existantes, le feature engineering \emph{cr\'ee} de nouvelles variables qui capturent des motifs que les donn\'ees brutes ne repr\'esentent pas explicitement.

\begin{datatip}
Une variable bien construite peut valoir plus qu'un mod\`ele complexe. Une simple r\'egression lin\'eaire avec les bonnes variables surpasse souvent un r\'eseau profond avec des variables brutes.
\end{datatip}

% ============================================================
\section{Variables polynomiales et d'interaction}

\begin{minted}{python}
import pandas as pd
import numpy as np
from sklearn.preprocessing import PolynomialFeatures

# Melbourne Housing
df = pd.read_csv("melb_data.csv")
df = df.dropna(subset=["Price", "Rooms", "Distance", "Landsize"])

# Variables polynomiales : capturer des relations non lin\'eaires
poly = PolynomialFeatures(degree=2, interaction_only=False,
                          include_bias=False)
features = df[["Rooms", "Distance"]]
poly_features = poly.fit_transform(features)
poly_names = poly.get_feature_names_out(["Rooms", "Distance"])

poly_df = pd.DataFrame(poly_features, columns=poly_names)
print(poly_df.head())
print(f"Original features: {features.shape[1]}")
print(f"Polynomial features: {poly_df.shape[1]}")
\end{minted}

\begin{minted}{python}
# Interactions seules : pas de termes au carr\'e, seulement les produits crois\'es
poly_inter = PolynomialFeatures(degree=2, interaction_only=True,
                                 include_bias=False)
inter_features = poly_inter.fit_transform(features)
inter_names = poly_inter.get_feature_names_out(["Rooms", "Distance"])
print(f"Interaction features: {inter_names}")
\end{minted}

\begin{warningbox}
Les variables polynomiales croissent combinatoirement : $n$ variables au degr\'e $d$ produisent $\binom{n+d}{d} - 1$ variables. Avec 20 variables au degr\'e 3, on obtient 1770 variables. \`A utiliser s\'electivement, pas sur tout le jeu de variables.
\end{warningbox}

% ============================================================
\section{Transformations math\'ematiques}

\begin{minted}{python}
# Variables de ratio : souvent plus informatives que les valeurs brutes
df["price_per_room"] = df["Price"] / df["Rooms"]
df["price_per_sqm"] = df["Price"] / df["Landsize"].replace(0, np.nan)
df["rooms_per_bathroom"] = df["Rooms"] / df["Bathroom"].replace(0, 1)

print(df[["Price", "Rooms", "price_per_room",
          "price_per_sqm"]].describe())
\end{minted}

\begin{minted}{python}
# Transformation log pour les variables asym\'etriques
df["log_price"] = np.log1p(df["Price"])
df["log_landsize"] = np.log1p(df["Landsize"])

# Racine carr\'ee : plus douce que log
df["sqrt_distance"] = np.sqrt(df["Distance"])

# Comparer la corr\'elation avec la cible
for col in ["Distance", "sqrt_distance", "log_price"]:
    if col in df.columns and df[col].notnull().sum() > 0:
        corr = df[col].corr(df["Price"])
        print(f"Corr({col}, Price) = {corr:.3f}")
\end{minted}

\begin{preprocesstip}
Les variables de ratio encodent de la \emph{connaissance m\'etier} en un seul nombre. << Prix au m\`etre carr\'e >> est une m\'etrique immobili\`ere standard qui capture mieux la densit\'e de valeur que prix ou surface seuls.
\end{preprocesstip}

% ============================================================
\section{Variables de date et d'heure}

\begin{minted}{python}
# Melbourne Housing : extraire des composantes de date
df["Date"] = pd.to_datetime(df["Date"], dayfirst=True)

df["sale_year"] = df["Date"].dt.year
df["sale_month"] = df["Date"].dt.month
df["sale_dayofweek"] = df["Date"].dt.dayofweek  # 0 = lundi
df["sale_quarter"] = df["Date"].dt.quarter
df["sale_is_weekend"] = df["Date"].dt.dayofweek.isin([5, 6]).astype(int)

print(df[["Date", "sale_year", "sale_month",
          "sale_dayofweek", "sale_quarter"]].head())
\end{minted}

\begin{minted}{python}
# Encodage cyclique du mois (pour que d\'ecembre et janvier soient proches)
df["month_sin"] = np.sin(2 * np.pi * df["sale_month"] / 12)
df["month_cos"] = np.cos(2 * np.pi * df["sale_month"] / 12)

# Jours depuis une date de r\'ef\'erence
reference = pd.Timestamp("2016-01-01")
df["days_since_ref"] = (df["Date"] - reference).dt.days

print(df[["sale_month", "month_sin", "month_cos"]].head(12))
\end{minted}

% ============================================================
\section{Variables d'agr\'egation}

\begin{minted}{python}
# Statistiques au niveau du quartier
suburb_stats = df.groupby("Suburb")["Price"].agg(
    suburb_mean_price="mean",
    suburb_median_price="median",
    suburb_price_std="std",
    suburb_n_sales="count"
).reset_index()

df = df.merge(suburb_stats, on="Suburb", how="left")

# Comment ce bien se compare-t-il \`a la moyenne du quartier ?
df["price_vs_suburb"] = df["Price"] / df["suburb_mean_price"]

print(df[["Suburb", "Price", "suburb_mean_price",
          "price_vs_suburb"]].head(10))
\end{minted}

\begin{minted}{python}
# Gapminder : variables par pays dans le temps
gap = pd.read_csv("gapminder.csv")

# Taux de croissance ann\'ee sur ann\'ee
gap = gap.sort_values(["country", "year"])
gap["gdp_growth"] = gap.groupby("country")["gdpPercap"].pct_change()
gap["pop_growth"] = gap.groupby("country")["pop"].pct_change()

# Moyenne mobile (tendance liss\'ee)
gap["lifeExp_rolling3"] = (gap.groupby("country")["lifeExp"]
                               .transform(lambda x: x.rolling(3).mean()))

print(gap[gap["country"] == "Benin"][
    ["year", "lifeExp", "lifeExp_rolling3", "gdp_growth"]].tail(6))
\end{minted}

% ============================================================
\section{Variables d\'eriv\'ees du texte}

\begin{minted}{python}
# Extraire des variables de colonnes texte sans NLP complet
# Sur les noms de quartiers et adresses Melbourne

# Longueur de l'adresse
df["address_length"] = df["Address"].str.len()
df["address_word_count"] = df["Address"].str.split().str.len()

# Contient des mots-cl\'es sp\'ecifiques
df["is_street"] = df["Address"].str.contains("St$", regex=True).astype(int)
df["is_road"] = df["Address"].str.contains("Rd$", regex=True).astype(int)
df["is_avenue"] = df["Address"].str.contains("Av", regex=False).astype(int)

print(df[["Address", "address_length", "is_street",
          "is_road"]].head(10))
\end{minted}

% ============================================================
\section{Variables m\'etier}

\begin{minted}{python}
# Connaissance m\'etier immobilier
df["total_rooms"] = df["Rooms"] + df["Bathroom"] + df["Car"]
df["has_multiple_bathrooms"] = (df["Bathroom"] > 1).astype(int)
df["is_new_build"] = (df["YearBuilt"] > 2010).astype(int)
df["building_age"] = df["sale_year"] - df["YearBuilt"]
df["is_close_to_cbd"] = (df["Distance"] < 10).astype(int)

# Indicateurs de type de bien
df["is_house"] = (df["Type"] == "h").astype(int)
df["is_unit"] = (df["Type"] == "u").astype(int)

# Interaction : grande maison proche du centre
df["large_close"] = df["is_house"] * df["is_close_to_cbd"] * df["Rooms"]

print(df[["Suburb", "Type", "Rooms", "Distance",
          "building_age", "large_close"]].head(10))
\end{minted}

\begin{datatip}
Les meilleures variables viennent de la compr\'ehension du domaine. Passez du temps avec les experts m\'etier. Demandez : << Quels facteurs les professionnels du domaine consid\`erent-ils dans leurs d\'ecisions ? >> Encodez ces facteurs en variables.
\end{datatip}

% ============================================================
\section{S\'election de variables apr\`es engineering}

\begin{minted}{python}
from sklearn.feature_selection import mutual_info_regression

# S\'electionner les variables num\'eriques
feature_cols = ["Rooms", "Distance", "Landsize", "building_age",
                "price_per_room", "suburb_mean_price", "large_close",
                "days_since_ref", "Bathroom", "Car"]
target = "Price"

# Supprimer les NaN pour cette analyse
subset = df[feature_cols + [target]].dropna()

mi = mutual_info_regression(subset[feature_cols], subset[target],
                            random_state=42)
mi_series = pd.Series(mi, index=feature_cols).sort_values(ascending=False)

print("Mutual information with Price:")
print(mi_series)
\end{minted}

\begin{minted}{python}
# Matrice de corr\'elation des variables construites
import seaborn as sns
import matplotlib.pyplot as plt

corr = subset[feature_cols].corr()
fig, ax = plt.subplots(figsize=(10, 8))
sns.heatmap(corr, annot=True, fmt=".2f", cmap="coolwarm",
            center=0, ax=ax)
plt.title("Feature correlation matrix")
plt.tight_layout()
plt.savefig("feature_correlation.png", dpi=150)
plt.show()
\end{minted}

% ============================================================
\section{Exercices}

\begin{exercise}
\begin{enumerate}
  \item Pour Melbourne Housing, cr\'eez au moins 5 nouvelles variables avec connaissance m\'etier. Classez-les par information mutuelle avec Price.
  \item Cr\'eez un encodage cyclique du mois de vente. V\'erifiez que d\'ecembre et janvier sont proches dans l'espace encod\'e en calculant leur distance euclidienne.
  \item Avec Gapminder, cr\'eez des taux de croissance ann\'ee sur ann\'ee pour GDP per capita et esp\'erance de vie. Quel pays a connu la plus forte croissance de PIB sur une p\'eriode ?
  \item Construisez des variables polynomiales (degr\'e 2) pour \texttt{Rooms}, \texttt{Distance}, et \texttt{Landsize}. Entra\^inez une r\'egression lin\'eaire avec et sans variables polynomiales et comparez le $R^2$.
  \item Cr\'eez une fonction d'engineering \texttt{engineer\_melb(df)} qui prend le dataset Melbourne brut et renvoie un DataFrame avec au moins 10 nouvelles variables, pr\^et pour la mod\'elisation.
\end{enumerate}
\end{exercise}

% ============================================================
\section{R\'esum\'e du chapitre}

\begin{itemize}
  \item Le feature engineering cr\'ee de nouvelles variables qui capturent des motifs invisibles dans les donn\'ees brutes.
  \item Les variables polynomiales et d'interaction capturent les relations non lin\'eaires et inter-variables.
  \item Les variables de ratio (prix par pi\`ece, prix au m\`etre carr\'e) encodent la connaissance m\'etier de fa\c{c}on compacte.
  \item Les variables de date (ann\'ee, mois, jour de semaine, encodage cyclique) d\'everrouillent les motifs temporels.
  \item Les variables d'agr\'egation (moyennes de groupe, moyennes mobiles) ajoutent du contexte.
  \item La connaissance m\'etier est la source la plus pr\'ecieuse d'id\'ees de variables --- consultez les experts.
  \item Apr\`es l'engineering, utilisez l'information mutuelle ou la corr\'elation pour s\'electionner les variables les plus pr\'edictives.
\end{itemize}
